\chapter{Android and Apple Platform Evidence}
\label{ch:android-macos}
\label{ch:android-ndk}
\index{Android}\index{Android emulator}\index{macOS}\index{iOS}\index{tvOS}

Platform support is not one fact.  A source tree may contain platform branches, cross-compile,
run tests in an emulator, package an application, draw a frame, and survive real hardware---or
stop at any rung in between.  Android retains recorded emulator evidence.  Alpha.1 contains
declared native macOS, iOS device-link and iOS simulator-launch routes, while tvOS remains outside
the implemented host.  Workflow source establishes those routes; without the tag's Actions run
records this edition does not promote each route to observed success.

Target OS, selected platform implementation and renderer are separate.  The Apple and Android
routes described here use the SDL3 platform implementation unless configured otherwise; iOS is a
target OS, not a fifth \texttt{CNA\_PLATFORM} value.  The iOS smoke selects
\texttt{SDL\_RENDERER}; native \texttt{METAL} remains macOS-only at the tag.

\section{Android application structure}

The Devices demo contains a complete Gradle and NDK application rather than a sketch.  Its 39
files include a Gradle wrapper, manifest, \texttt{SDLActivity} subclass, SDL Java glue, launcher
resources, and native build.  The JNI-side CMake file uses
\texttt{add\_subdirectory(<cna-root>)}, so producing the APK compiles the CNA library as part of
the application.  The recorded configuration uses compile SDK~35, NDK~30, minimum SDK~24, and
only the \texttt{arm64-v8a} ABI.

Android-specific production code is concentrated rather than scattered. The Devices module
contains about 2,090 lines of NDK sensor implementation using \texttt{ASensorManager}, an
\texttt{ALooper}, event queues and rollback. Each \texttt{AndroidSensorBridge} can own a dedicated
polling thread; the Motion backend aggregates six bridge instances, so the single
\texttt{std::thread} construction site can produce several concurrent workers. CNA itself uses
no JNI. Its only Android-specific link adds \texttt{android} and \texttt{log} to the Devices
module; the root CMake file contains no Android branch. Portable sensor mathematics stays outside
\texttt{\_\_ANDROID\_\_} guards so native tests can exercise it, while the I/O boundary is
platform-gated.

\section{Recorded Android execution}

Two emulator campaigns are recorded. On 4~July~2026 an NDK-built \texttt{CnaTests} executable
was pushed to a Medium Phone AVD. All 230 Net/Gamer/ENet/Packet cases and 1,831 other cases passed
before an unrelated \texttt{TitleContainer} test reached the no-Activity crash described below;
surviving build-system repairs corroborate the run. On 5~July the Devices application produced a
7.3\,MB APK, installed and started it, captured the emulator screen, and injected sensor data
through the emulator console. The commit record says the real UI and its sensor-event flash were
visible in that screenshot. That injection exercised SDL's Accelerometer/Gyroscope delivery;
Compass/Motion's CNA-owned NDK bridge and six-source fusion were not separately driven.

Those are meaningful execution results, but the repository contains no APK, AAB, logcat log,
Android screenshot, or machine-readable test result from them.  The strongest honest wording is
therefore ``recorded emulator execution,'' not a continuously reproducible Android gate.  No CNA
binary has been evidenced on physical Android hardware at the audited revision.

Android graphics is a separate and still open boundary, but ``never attempted'' is too broad. The
July Devices APK's nested CNA build used Android's default \texttt{SDL\_RENDERER}, and its commit
records screenshot-confirmed UI presentation. That is historical renderer execution, not a
retained pixel oracle: no screenshot survives in the repository, no graphics-specific assertion
was made, and the later adaptation was not rebuilt at the audited pin. WebGPU is the only renderer
containing an explicit Android native-window surface path, and it has not run there. The strongest
claim is thus recorded pre-adaptation SDL-renderer presentation on an AVD, not current Android
graphics verification and not evidence for OpenGL or WebGPU.

\section{Lifecycle coverage is explicit but partial}

The platform event contract carries will-enter-background, did-enter-foreground, low-memory and
terminating events.  \cnaclass{Game} suspends between the first two, dispatches low-memory to
content cleanup and exits on termination.  That is source-level lifecycle policy plus portable
tests; it does not prove every mobile host delivers the events in the expected order or that a
game's GPU resources survive a background transition.

This is best expressed as a state-machine gap.  Launch-to-visible behavior covers only the
straight path.  A useful mobile test must drive background, foreground, low-memory, and process
recreation transitions, then verify saved state and graphics-resource reconstruction.  An
emulator can provide most of this evidence before physical-device testing adds drivers, thermal
behavior, and vendor-specific lifecycle timing.

\section{Assets expose a deeper integration gap}

\texttt{TitleContainer} first uses \texttt{std::filesystem}; on Android it can fall back to
\texttt{SDL\_LoadFile} and wrap the bytes in a memory stream.  Historical execution found that
fallback could segfault without a real Activity/JNI context.  More importantly, no production
consumer calls \texttt{TitleContainer}: \texttt{ContentManager} builds
\texttt{Content/<asset>} directly from the current working directory.

The fallback is therefore both fragile and unreachable from the main content path.  The same
CWD-relative content root is a latent issue for Android packages and macOS application bundles.
A complete mobile asset contract should choose one route, base it on the package location or
platform asset API, and test it through \texttt{ContentManager} rather than through an isolated
helper.

\section{macOS: native routes are declared}

The Metal renderer keeps Objective-C\texttt{++} behind a pure-C\texttt{++} PIMPL boundary.
Objective-C\texttt{++}, Metal, QuartzCore, and Foundation are
enabled only when the selected identity is \texttt{METAL}.  This is a sound portability seam,
but source isolation is not native validation.

The \texttt{Metal} CTest prefix is much broader than three registrations. Portable Metal policy
gtests are discovered individually and also collected in \texttt{Metal\_PortableHelpers}; a
recorded Linux HEADLESS run consequently reports 207/207 for \texttt{ctest -R '\^Metal'} while
executing no Objective-C\texttt{++}, Apple framework, shader compiler, or GPU command. Only the
separate \texttt{Metal\_Smoke} and \texttt{Metal\_Capabilities} registrations require a
\texttt{METAL} configuration. The renderer identity of the build must accompany every quoted
result.

There is genuine historical Apple evidence: a macOS~14/Xcode~15.4 workflow on the original
feature tree ran 143 tests, passing 136 and failing seven.  The later adaptation changed
interfaces, ownership, transfer policy, registration, and supported behavior.  That run is
valuable history, not evidence for alpha.1.

At the tag, the dedicated automatic Metal workflow configures \texttt{METAL} on macOS~14, builds
the native target, enables Metal shader validation/debug-layer variables and runs the Metal CTest
prefix plus a renderer-definition control.  The separate Apple workflow builds an
\texttt{SDL\_RENDERER} macOS suite and a self-contained \texttt{.app}, checks Mach-O dependency
paths and launches it for a success marker.  These definitions are materially stronger than the
old stale-filter route, but this edition did not retrieve workflow-run artifacts and therefore
records declared gates rather than a new pixel/Retina/frame-pacing observation.

\section{iOS: final-link and simulator route; tvOS remains a non-claim}

The tag now contains an iOS toolchain, application metadata and a minimal UIKit/SDL-main/Game
smoke application.  Its workflow matrix final-links \texttt{SDL\_RENDERER} application bundles
for arm64 device and arm64 simulator, inspects Mach-O platform load commands, verifies the SDL
entry symbols and absence of build-machine dynamic dependencies, then ad-hoc signs, installs and
launches the simulator bundle until a one-frame success marker appears.  Tests/examples and Net
are deliberately off; no physical device is run.  This is an executable route in source, not
evidence that the tag's matrix ran successfully unless its Actions record is also retained.

tvOS has no corresponding toolchain/application/workflow route and remains an explicit
non-claim.  Native \texttt{METAL} is also not an iOS claim: the available smoke uses
\texttt{SDL\_RENDERER} over the SDL3 platform implementation.

The next durable milestone is intentionally narrower than ``mobile support'': archive a fresh
Android graphics build and emulator frame with logs; retain alpha.1 Metal and Apple workflow
artifacts; add a physical iOS-device result; then design, rather than infer, a tvOS host.  Each step should state
which rung it proves.  That makes platform progress cumulative without allowing yesterday's
partial result to masquerade as today's complete one.
